\chapter{Child-Order Placement}\label{ch:mx:child-order-placement}

An algorithm must buy 800 shares in the next minute. It can join the queue at the bid and hope, or pay the spread now. Waiting saves half a spread
if it fills and costs more than that if the price walks away while it waits. This chapter states the choice as a problem, measures the obvious
policies in the simulated market, solves a Markov model of the best quotes by dynamic programming, lets a tabular learner find the same answer, and
looks for the queue imbalance above which crossing at once is cheaper.

\section{The placement problem}

\begin{definition}[Order placement problem]\label{def:mx:child-order-placement:problem}
The \emph{order placement problem}\index{order placement problem} is the choice, for a \omterm{def:mx:the-almgren-chriss-framework:parent}{child order} that must be done by a deadline, of where and when
to post it (at the best quote, behind it, inside the spread, or across it as a marketable order) and of how to revise it, so as to minimise the
expected cost against the mid at the decision, penalised by its dispersion.
\end{definition}

\begin{definition}[Non-execution risk, clean-up trade]\label{def:mx:child-order-placement:risk}
\emph{Non-execution risk}\index{non-execution risk} is the chance that a passive order is not filled by its deadline, and the cost of what must then
be done. The \emph{clean-up trade}\index{clean-up trade} is that remedy: at the deadline the unfilled remainder is cancelled and sent as a marketable
order, at whatever the price has become.
\end{definition}

A \omterm{def:mx:the-almgren-chriss-framework:parent}{child order} (chapter~14) inherits a quantity and a deadline from the schedule of chapter~16. Its cost is one of three things. Crossing at once pays
half the spread plus whatever the order walks through beyond the best quote. Joining the bid and filling earns half the spread; its fill probability is
the queue's business (chapter~6's \omterm{def:mx:tick-size-queues-and-priority:queuevalue}{queue value}, One Quant Book 7, chapter~18's passive fill probability). Joining and not filling adds the \omterm{def:mx:child-order-placement:risk}{clean-up
trade} after the price has moved away, the adverse selection (One Quant Book 1, chapter~1) of the orders that do not fill. Harris and Hasbrouck (1996)
measured the trade-off on NYSE SuperDOT orders and found that \omterm{def:mx:the-limit-order-book:orders}{limit orders} placed at or better than the quote beat \omterm{def:mx:the-limit-order-book:orders}{market orders} even after a penalty
for the orders that did not execute.

The simulated market of this chapter is \texttt{firm.agentmkt} with calmer liquidity providers than chapter~16's: on a 30-minute session the spread is
one tick 97.7\% of the time, the best quotes hold 14.6 lots (of 100 shares) on average, the mid moves once every 9.4 seconds and 137 shares trade a
second, so the 800-share slice is about a tenth of a minute's volume.

\section{Passive against aggressive}

Five policies work the same nine slices a session (buys and sells in turn, two minutes apart), on 24 sessions with the same order flow for every
policy:
\begin{itemize}
\item \emph{cross}: a marketable order for the 800 shares at once;
\item \emph{join}: a \omterm{def:mx:the-limit-order-book:orders}{limit order} at the bid, left alone, cleaned up at the deadline;
\item \emph{behind}: the same one tick below the bid;
\item \emph{reprice}: join, and follow the bid when it moves away;
\item \emph{imbalance}: cross at once if the queue imbalance on the order's side, $(B-A)/(B+A)$ with $B$ the size at the bid and $A$ at the ask for a
buy (One Quant Book 7, chapter~8), exceeds 0.5, a common rule of thumb; otherwise reprice.
\end{itemize}

\begin{definition}[Repricing rule]\label{def:mx:child-order-placement:reprice}
A \emph{repricing rule}\index{repricing rule} says when a resting \omterm{def:mx:the-almgren-chriss-framework:parent}{child order} is cancelled and replaced at a new price: here, whenever the best
quote on its side moves away from it, at the cost of its place in the queue.
\end{definition}

\begin{table}[ht]
\centering
\small
\begin{tabular}{@{}lrrrrr@{}}
\toprule
policy & cost (s.e.) & s.d. & against crossing (s.e.) & filled passively & cleaned up \\
\midrule
cross & 0.70 (0.04) & 0.53 & --- & 0\% & 0\% \\
join & $-0.24$ (0.06) & 0.89 & $-0.94$ (0.07) & 91.0\% & 9.0\% \\
behind & 0.76 (0.11) & 1.59 & 0.06 (0.11) & 16.7\% & 83.3\% \\
reprice & $-0.22$ (0.05) & 0.72 & $-0.92$ (0.06) & 97.3\% & 2.7\% \\
imbalance & $-0.19$ (0.05) & 0.78 & $-0.89$ (0.06) & 90.8\% & 2.7\% \\
\bottomrule
\end{tabular}
\caption{Working an 800-share \omterm{def:mx:the-almgren-chriss-framework:parent}{child order} within 60 seconds in the simulated market: 216 slices, cost per share in ticks against the mid at the
decision (positive: paid more), its standard deviation, the paired difference to crossing, and the shares filled passively or by the \omterm{def:mx:child-order-placement:risk}{clean-up trade}.
Data: \texttt{mx\_place.sim\_study}.}\label{tab:mx:child-order-placement:policies}
\end{table}

Joining the bid costs 0.94 ticks a share less than crossing, almost the whole spread: the queue fills 91\% of the order in the minute. Crossing is
expensive not for its half-tick but for the 3.2 lots of the second level it walks into (0.70 ticks on average). Stepping one tick behind almost never
fills (16.7\%) and leaves 83\% to the \omterm{def:mx:child-order-placement:risk}{clean-up trade} after the price has moved: it is the worst policy and the most dispersed. Repricing trades a few
hundredths of a tick of mean for a smaller spread of outcomes (0.72 against 0.89): the order it saves from the clean-up is the one the price ran away
from. The imbalance rule crossed at the start of 6.5\% of the slices and paid for it: 0.03 ticks worse than repricing on average.

\section{Queue management and repricing}

Two facts drive a resting order: where it is in the queue and whether the queue will be there when its turn comes. Its place is known from the
feed (\texttt{firm.exchsim}'s \texttt{working()} counts the shares ahead in the agent's own view of the book); the queue's future is what the imbalance
predicts: a thin ask facing a deep bid tends to be consumed first, and the price ticks up away from a resting buy. Lehalle and Mounjid (2017) built
limit-order strategies that monitor this imbalance to reduce adverse selection.

\begin{definition}[Queue-aware placement]\label{def:mx:child-order-placement:qaware}
\emph{Queue-aware placement}\index{queue-aware placement} conditions each decision (wait, reprice, cross) on the order's position in its queue and on
the sizes of the queues on both sides, not only on the prices and the time left.
\end{definition}

A \omterm{def:mx:child-order-placement:reprice}{repricing rule} has two failure modes. Too eager, it gives up queue position for nothing when the quote flickers; too lazy, it leaves the order behind
the market, where \emph{behind} ended. Real engines add a minimum resting time, a tolerance of a tick before chasing, and a limit on how far to chase
before crossing; on this market the simple rule already captures most of the value.

\section{Placement as a Markov decision problem}

Model the best quotes as queues (the birth--death processes of One Quant Book 4, chapter~8, and Huang, Lehalle and Rosenbaum's queue-reactive
model, which makes the rates depend on the queues' sizes). A buy of $r$ lots rests behind $n$ lots at the bid; the ask holds $a$ lots. In a step
$dt$, a market sell arrives with probability $\mu\,dt$ and takes one lot from the front (one of ours when $n=0$, earning half a tick against the mid),
a lot ahead of us cancels with probability $\theta n\,dt$, the ask loses a lot with probability $(\mu+\theta a)\,dt$ and gains one with $\lambda\,dt$.
When the ask empties, the price moves up a tick and the remaining lots start again from fresh queues, a tick worse. At each step the order may
cross: $r$ lots against $a$ at the ask cost half a tick each for the first $a$ and 1.5 for the next level. Dynamic programming (One Quant Book 4,
chapter~9) gives the value $V_k(r,n,a)$ with $k$ steps left:
\[
V_k=\min\bigl(C(r,a),\ \E[V_{k-1}\mid r,n,a]\bigr),\qquad V_0=C(r,a).
\]

\omcode{code/firm/placement/firm_placement.py}{147}{159}{The backward induction of the placement problem: for each number of lots left, the value of waiting one step against crossing now, on the grid of lots ahead and lots at the ask.}\label{lst:mx:child-order-placement:dp}

Fitted on the feed of the calibration session (\texttt{firm\_placement.calibrate}), $\lambda=0.99$ lots a second join each best quote, each lot at the
best cancels at $\theta=0.021$ a second, \omterm{def:mx:the-limit-order-book:orders}{market orders} take $\mu=0.69$ lots a second from each side, and the second level holds 3.2 lots. For the
800-share slice the plan never crosses at the start: on the model, joining is worth $-0.40$ ticks a share against 0.66 for crossing. Crossing becomes
right only for small, urgent orders. For a single lot the plan crosses when the imbalance exceeds a threshold that falls as the deadline nears
(\cref{fig:mx:child-order-placement:threshold}): 0.86 with fifteen seconds left, 0.79 with five, 0.66 with three. For eight lots it never crosses
at the start, with fifteen seconds or sixty.

\begin{omfigure}
\begin{tikzpicture}
\begin{axis}[width=10cm,height=5.4cm,xlabel={\small seconds left},ylabel={\small imbalance threshold},
  tick label style={font=\footnotesize},grid=major,grid style={gray!25},xmin=0,xmax=21,ymin=-0.6,ymax=1,table/col sep=comma]
\addplot[omDef,very thick,mark=*,mark size=1.4pt] table[x=seconds,y=cut]{figdata/microstructure/17-child-order-placement/threshold.csv};
\addplot[gray,dashed,domain=0:21] {0.5};
\node[gray,font=\scriptsize,anchor=south east] at (axis cs:20.5,0.5) {rule of thumb};
\end{axis}
\end{tikzpicture}

\omcaption{The queue imbalance above which the plan crosses a single lot at once, against the time left: the single-threshold rule that best matches
the dynamic-programming decision over the model's fresh queues. Data: \texttt{mx\_place.threshold\_curve}.}\label{fig:mx:child-order-placement:threshold}
\end{omfigure}

\omcode{code/firm/placement/firm_placement.py}{330}{341}{The plan as a policy of the slice executor: the lots left, the lots ahead (or a fresh place at the back after the price moved) and the lots at the far quote, looked up in the dynamic-programming table.}\label{lst:mx:child-order-placement:policy}

Run in the simulated market, the plan's decisions coincide with repricing on every slice: it never found crossing worthwhile before the deadline.
The model is too optimistic about waiting, though: it expects $-0.40$ ticks a share where the market gives $-0.22$, because its price moves only
when a queue empties while the simulated fundamentalists push the price toward a value the model does not see.

Is the model's threshold real? One lot with five seconds left, 1\,824 slices, crossing and repricing paired on the same flow:
\cref{fig:mx:child-order-placement:urgent} shows the difference by imbalance. Waiting is 0.26 ticks cheaper when the imbalance is below $-0.5$; the
saving falls to a few hundredths above zero and is indistinguishable from nothing above 0.5 ($+0.03\pm0.07$ between 0.5 and 0.75, $-0.11\pm0.08$
above). The simulated market agrees that the advantage of waiting vanishes where the model says it does, but at no imbalance is crossing significantly
cheaper.

\begin{omfigure}
\begin{tikzpicture}
\begin{axis}[width=10cm,height=5.4cm,xlabel={\small queue imbalance on the order's side at the start},
  ylabel={\small join minus cross (ticks)},tick label style={font=\footnotesize},grid=major,grid style={gray!25},
  xmin=-1,xmax=1,ymin=-0.5,ymax=0.3,table/col sep=comma]
\addplot[omDef,very thick,mark=*,error bars/.cd,y dir=both,y explicit] table[x=mid,y=diff,y error=se]{figdata/microstructure/17-child-order-placement/urgent.csv};
\addplot[black,dashed,domain=-1:1] {0};
\end{axis}
\end{tikzpicture}

\omcaption{One lot with five seconds left: the cost of joining the bid (with repricing and the \omterm{def:mx:child-order-placement:risk}{clean-up trade}) minus the cost of crossing at once, by
queue imbalance at the start, with standard errors; 1\,824 paired slices. Data: \texttt{mx\_place.urgent\_study}.}\label{fig:mx:child-order-placement:urgent}
\end{omfigure}

Cont and Kukanov (2017) posed the placement problem across several venues, with fees and rebates, as a convex optimisation of how many shares to post
where and how many to send as \omterm{def:mx:the-limit-order-book:orders}{market orders}; chapter~18 takes up the routing half.

\section{Learning a placement policy}

The plan needed the model's rates. A learner can find a policy from experience alone: tabular Q-learning (One Quant Book 12, chapter~17, develops
reinforcement learning) estimates the cost-to-go $Q(s,a)$ of each state and action from simulated episodes and acts greedily. Nevmyvaka, Feng and
Kearns (2006) applied reinforcement learning to trade execution on NASDAQ order-book data, with a state of time left, inventory and market variables.
Here \texttt{firm\_placement.qlearn} runs 200\,000 episodes of the Markov book with the state (lots left, lots ahead, lots at the ask, five-second
bucket of time) and the actions wait or cross. Its greedy policy costs $-0.39$ ticks a share on 40\,000 fresh episodes against $-0.40$ for the plan:
close, in a state space of which it visited only 4.0\%. The learner reached the plan's answer (join and wait) without being told the rates; it could
not do better, and in a larger state space it would need far more episodes, or a function approximator, which is Book~12's subject.

\section{Tutorial: working a child order}\label{tut:mx:child-order-placement:tut}

\begin{tutorial}
\textbf{Goal.} Compare placement policies for a one-minute slice in the simulated market, solve the placement problem by dynamic programming on a
fitted Markov book, and learn it with Q-learning.
\textbf{End state:} \cref{tab:mx:child-order-placement:policies}, \cref{fig:mx:child-order-placement:threshold,fig:mx:child-order-placement:urgent}
and the numbers of sections 4 and 5.
\begin{tutsteps}
\item \textbf{Market.} \texttt{mx\_place.MARKET}, \texttt{market\_stats()}; \texttt{firm\_placement.calibrate(res)}.
\item \textbf{Policies.} \texttt{SliceExecutor(slices, policy)} with \texttt{Cross}, \texttt{Post(offset, reprice)}, \texttt{Imbalance},
\texttt{DPPolicy}; \texttt{sim\_study()}.
\item \textbf{Plan.} \texttt{solve(model, r\_max, horizon\_s)}, \texttt{imbalance\_threshold(plan, r, k)}; \texttt{dp\_thresholds()},
\texttt{threshold\_curve()}.
\item \textbf{Check.} \texttt{urgent\_study()}: one lot, five seconds, paired by imbalance.
\item \textbf{Learn.} \texttt{qlearn(model, \dots)}, \texttt{q\_policy(q)}, \texttt{simulate(model, policy, \dots)}; \texttt{model\_study()}; draw
with \texttt{fig\_place.py}.
\end{tutsteps}
\textbf{What to change next.} Give the fundamentalists more weight (stronger short-term drift) and find the imbalance at which crossing becomes
cheaper in the simulated market; add a mid-price drift to the Markov book and see the plan's threshold fall; replace the Q-table by a small network.
\end{tutorial}

\section{Build: placement}\label{bld:mx:child-order-placement:placement}

\begin{build}
\bfield{Purpose} The child-order layer under chapter~16's schedules and chapter~28's \omterm{def:mx:benchmark-algorithms:algo}{execution algorithm}: policies, a fitted book model, its optimal
plan and a learned baseline.
\bfield{Interface} \texttt{BookModel(lam, theta, mu, fresh, depth2)}, \texttt{calibrate(res)}, \texttt{cross\_cost(r, a, depth2)},
\texttt{solve(model, r\_max, horizon\_s, dt, n\_max, a\_max)} (\texttt{Plan.decide}, \texttt{.value}, \texttt{.fresh\_value}),
\texttt{imbalance\_threshold(plan, r, k)}, \texttt{simulate(model, policy, r, horizon\_s, \dots)}, \texttt{qlearn(\dots)}, \texttt{q\_policy(q)},
\texttt{SliceExecutor(slices, policy, check\_s)}, \texttt{Cross}, \texttt{Post}, \texttt{Imbalance}, \texttt{DPPolicy}.
\bfield{Rules} Lots of 100 shares; costs in ticks against the mid at the decision, per share or summed over lots; the clean-up at the deadline is
part of every passive policy; sells are the mirror image.
\bfield{Acceptance tests} \texttt{code/firm/placement/tests/}: crossing walks the book; the plan's value equals its own Monte Carlo cost and is no worse
than always or never crossing; the crossing region grows with the lots ahead and shrinks with the ask; Q-learning approaches the plan; the executor
completes cross and repricing slices.
\bfield{Stretch} Queue-reactive rates; a drift signal in the state; several venues (chapter~18).
\end{build}

\begin{omsources}
\item L.~Harris and J.~Hasbrouck, ``Market vs.\ limit orders: the SuperDOT evidence on order submission strategy'', \emph{Journal of Financial and
Quantitative Analysis} 31(2), 1996.
\item Y.~Nevmyvaka, Y.~Feng and M.~Kearns, ``Reinforcement learning for optimized trade execution'', \emph{Proceedings of the 23rd International
Conference on Machine Learning}, 2006.
\item W.~Huang, C.-A.~Lehalle and M.~Rosenbaum, ``Simulating and analyzing order book data: the queue-reactive model'', \emph{Journal of the American
Statistical Association} 110(509), 2015.
\item R.~Cont and A.~Kukanov, ``Optimal order placement in limit order markets'', \emph{Quantitative Finance} 17(1), 2017.
\item C.-A.~Lehalle and O.~Mounjid, ``Limit order strategic placement with adverse selection risk and the role of latency'', \emph{Market
Microstructure and Liquidity} 3(1), 2017.
\end{omsources}

\section{Exercises}

\begin{exercise}[$\star$]\label{exo:mx:child-order-placement:1}
A buy of 8 lots crosses an ask of 5 lots with 3.2 lots at the next level. What does it cost per share against the mid, with a one-tick spread?
\end{exercise}

\begin{exercise}[$\star$]\label{exo:mx:child-order-placement:2}
With the fitted rates, how long on average does a lot at the back of a 15-lot bid queue wait for the queue ahead to clear, if the queue ahead only
shrinks (market sells, and cancellations at $\theta$ per lot ahead) and the price does not move?
\end{exercise}

\begin{exercise}[$\star$]\label{exo:mx:child-order-placement:3}
Why does stepping one tick behind the bid cost more than crossing, when it pays no spread when it fills?
\end{exercise}

\begin{exercise}[$\star\star$]\label{exo:mx:child-order-placement:4}
Why does the plan's threshold fall as the deadline nears?
\end{exercise}

\begin{exercise}[$\star\star$]\label{exo:mx:child-order-placement:5}
The model expects $-0.40$ ticks a share for joining and the simulated market gives $-0.22$. Name two things the model leaves out that explain the gap.
\end{exercise}

\begin{exercise}[$\star\star$]\label{exo:mx:child-order-placement:6}
Why does repricing lower the standard deviation of the cost more than its mean?
\end{exercise}

\begin{exercise}[$\star\star\star$]\label{exo:mx:child-order-placement:7}
\emph{Coding.} Solve the plan for one lot with five seconds left. In what share of the fresh-queue states (weighted by their law) does it cross at
once, with the fitted rates and with $\lambda$ halved (fewer sellers joining the ask, which then empties sooner)?
\end{exercise}

\begin{exercise}[$\star\star\star$]\label{exo:mx:child-order-placement:8}
\emph{Find the flaw.} ``Our passive fills cost minus half a tick against the mid at the fill, so passive execution earns half a tick on every share.''
\end{exercise}

\section{Problem: Join, Step Ahead or Cross?}

\begin{problem}[{Weekend problem --- join, step ahead or cross?}]\label{pb:mx:child-order-placement:1}
An algorithm's \omterm{def:mx:the-almgren-chriss-framework:parent}{child orders} cross the spread every time; the desk asks whether they should rest instead, and when crossing is right.

\medskip\noindent\textbf{Part I --- The problem.}
\begin{enumerate}
\item State the \omterm{def:mx:child-order-placement:problem}{order placement problem} and its cost.
\item Define \omterm{def:mx:child-order-placement:risk}{non-execution risk} and the \omterm{def:mx:child-order-placement:risk}{clean-up trade}.
\item Why does stepping ahead of the bid hardly exist in the simulated market?
\item Describe the simulated market's spread, queues and pace.
\end{enumerate}

\medskip\noindent\textbf{Part II --- Policies in the simulated market.}
\begin{enumerate}[resume]
\item Give the cost per share and its standard deviation for crossing, joining, stepping behind and repricing.
\item Why is crossing 800 shares more expensive than half a tick?
\item What share of the order does joining fill passively, and what share goes to the \omterm{def:mx:child-order-placement:risk}{clean-up trade}?
\item What did the imbalance rule of thumb cost, and how often did it cross?
\end{enumerate}

\medskip\noindent\textbf{Part III --- The plan.}
\begin{enumerate}[resume]
\item Write the Markov book and its transitions.
\item Write the backward induction.
\item Give the fitted rates.
\item What does the plan do with the 800-share slice, and what value does it expect?
\item Give the plan's imbalance threshold for one lot with fifteen, five and three seconds left.
\item \emph{State the named result}: the expected cost per share of each placement policy, and the queue-imbalance threshold above which crossing
at once is cheaper.
\item What does the simulated market say about that threshold for one lot with five seconds left?
\end{enumerate}

\medskip\noindent\textbf{Part IV --- Learning and the answer.}
\begin{enumerate}[resume]
\item How does tabular Q-learning estimate the policy, and on what state?
\item How close does it come to the plan, and how much of its table did it visit?
\item Why would it struggle on a real book?
\item What would you tell the desk?
\item In one sentence: when is crossing at once right?
\end{enumerate}
\end{problem}

\section{Interview questions}

\begin{interviewq}[$\star$ \iqroles{trader}]\label{iq:mx:child-order-placement:1}
You must buy 1\,000 shares in the next minute in a stock with a one-tick spread and deep queues. \omterm{def:mx:the-limit-order-book:orders}{Market order} or \omterm{def:mx:the-limit-order-book:orders}{limit order}?
\end{interviewq}

\begin{interviewq}[$\star\star$ \iqroles{researcher}]\label{iq:mx:child-order-placement:2}
How does queue imbalance enter the decision to cross, and why?
\end{interviewq}

\begin{interviewq}[$\star\star$ \iqroles{researcher}]\label{iq:mx:child-order-placement:3}
Set up child-order placement as a dynamic programme. What is the state?
\end{interviewq}

\begin{interviewq}[$\star\star$ \iqroles{developer}]\label{iq:mx:child-order-placement:4}
Your repricing logic cancels and replaces an order every time the bid flickers. What goes wrong, and what would you add?
\end{interviewq}

\begin{interviewq}[$\star\star$ \iqroles{mle}]\label{iq:mx:child-order-placement:5}
Would you use reinforcement learning for order placement? What would you compare it with, and how?
\end{interviewq}

\begin{interviewq}[$\star\star\star$ \iqroles{trader}]\label{iq:mx:child-order-placement:6}
Your passive fill rate is 90\% but your execution cost against arrival is worse than a colleague who crosses more. How can that be?
\end{interviewq}
